To test whether the feature-based models understate the attainable ceiling, ROCKET \cite{dempster2020}, a state-of-the-art time-series method that applies 500 random convolutional kernels to the raw 1440-minute series and fits a regularised logistic head on the resulting 1{,}000 features, was evaluated under both designs (Table~\ref{tab:e9}). Its subject-wise AUROC was 0.717 (95\% CI 0.624 to 0.799) on DEPRESJON, 0.865 on PSYKOSE and 0.506 on HYPERAKTIV, below the feature-based models in every cohort, and record-wise splitting inflated it by 0.114 (0.066 to 0.167), 0.040 (0.021 to 0.068) and 0.069 (0.041 to 0.092), with participant-level leaky estimates of 0.976, 0.980 and 0.657. Its calibration slopes under honest splitting were 0.36, 0.67 and 0.00. A stronger learner therefore did not raise the honest ceiling; it reproduced the same inflation.
